\documentclass{jarticle}
\usepackage[dvipdfmx]{graphicx}
\usepackage{amsmath}
\usepackage{amssymb}
\usepackage{bm}
\usepackage{float}
\usepackage{xcolor}

% 体裁（必要なら調整）
\topmargin -0.4mm
\oddsidemargin -0.4mm
\evensidemargin -0.4mm
\headheight 0mm
\headsep 0mm
\footskip 10mm
\textheight 247mm
\textwidth 170mm
\setlength{\parindent}{1zw}
\pagestyle{empty}

\makeatletter
\def \fgcaption{\def \@captype{figure} \caption}
\def \tbcaption{\def \@captype{table} \caption}
\makeatother

\def\figurename{Fig.}
\def\tablename{Table}
\def\topfraction{.95}
\def\bottomfraction{.95}
\def\textfraction{.05}
\def\floatpagefraction{.8}
\def\dbltopfraction{.95}
\def\dbltextfraction{.05}
\def\dblfloatpagefraction{.8}

\begin{document}

%==============================
%  ヘッダー
%==============================
\begin{center}
\begin{tabular}{|c|c|c|} \hline
  \parbox{50mm}{\vspace{1mm} \centering{\Large\bf 2026/04/15}} &
  \parbox{55mm}{\vspace{1mm} \centering{\Large\bf M1\quad 市川 詠祥}} &
  \parbox{50mm}{\vspace{1mm} \centering{\Large\bf No.~2}} \ [2mm] \hline

  \multicolumn{3}{|c|}{
    \parbox{155mm}{
      \centering\Large\bf \vspace{2mm}
       くさび形モデルにおける表面波FDTDシミュレーションの実装
    }
  } \ [3mm] \hline\hline

  \multicolumn{3}{|c|}{
    \parbox{155mm}{
      \large\vspace{2mm}

      {\Large\bf 今月の目標}
      \begin{itemize}
        \item くさび形のきずにおける表面波の解析
      \end{itemize}

      {\Large\bf 進捗の要点}
      \begin{itemize}
        \item くさび形FDTDシミュレーションコードの整備（軸定義・形状名の統一）
        \item くさび形きずにおける表面波シミュレーションプログラムの新規作成
        \item 縦方向応力（T1）・横方向応力（T3）の時空間マップ可視化プログラムの実装
      \end{itemize}

      {\Large\bf 今後の実施事項}
      \begin{itemize}
        \item くさび形きずにおける表面波の特性解析
        \item U字形モデルとの比較検討
      \end{itemize}

      {\Large\bf 投稿・発表予定}
      \begin{itemize}
        \item
      \end{itemize}
    }
  } \ [3mm] \hline
\end{tabular}
\end{center}

\vspace{6mm}
{\Large\bf 今回の内容}
\vspace{3mm}

\subsection*{1. 新入生向け発表（4/13）}
月曜日に新入生向けの研究紹介発表を行った．
先週末はその準備に充てていたため，シミュレーション関連の作業は月曜日以降から開始した．

\subsection*{2. コード・リポジトリの整理（4/13）}
GitHubリポジトリの管理構成を整備した．
\texttt{project2}（別PC由来のコード群）を\texttt{.gitignore}に追加し，GitHubへの同期対象を\texttt{project4}のみに限定した．
また，\texttt{README.md}をproject4用に書き直し，\texttt{project2/README.md}を分離して「別PCで作成されたプログラム群である」旨の注意書きを追記した．

\subsection*{3. 表面波FDTDシミュレーションプログラムの作成（4/14）}
くさび形きずにおける表面波を観測するためのFDTDシミュレーションプログラム\texttt{kusabi\_surface\_sim.py}を新規作成した．
このプログラムでは，縦方向応力$T_1$および横方向応力$T_3$を複数の観測点で時系列記録し，時空間マップとして保存できる構成とした．
また，探触子配置と計測点配置を確認するための可視化スクリプト（\texttt{visualize\_measurement\_points.py}等）も合わせて作成した．

\subsection*{4. 時空間マップ可視化プログラムの実装（4/15）}
論文のFig.~3.5に相当する図を再現するため，可視化プログラム\texttt{plot\_kusabi\_surface\_map.py}を新規作成した．
このプログラムは，シミュレーション結果（.npzファイル）を読み込み，
$T_1$（z方向の固定断面における縦方向応力の$x$-time マップ）と
$T_3$（表面計測点における横方向応力の$z$-time マップ）を
1枚の図にまとめて出力する．
あわせて\texttt{kusabi\_surface\_sim.py}のリファクタリングを行い，コードを整理した．

%===========================================

\end{document}
